% Introduction section for the ESIPS honours thesis.
% Intended to be pulled into main.tex via % Introduction section for the ESIPS honours thesis.
% Intended to be pulled into main.tex via % Introduction section for the ESIPS honours thesis.
% Intended to be pulled into main.tex via % Introduction section for the ESIPS honours thesis.
% Intended to be pulled into main.tex via \input{introduction} (or \include).
% Assumes an IEEEtran two-column document; contains no preamble.
% Section cross-references use \ref{} to auto-numbered labels. The one exception is
% the Conclusion (Section~VI), left as a literal numeral until that section is written
% and wired into main.tex; see the PLACEHOLDER note at the end of this file.

\section{Introduction}
\label{sec:intro}

A literature search asks two things of a retrieval system: firstly, to find papers that precisely match a set of stated constraints; secondly, to find papers that are about what the user means. These map onto two retrieval paradigms that have proven difficult to reconcile. Structured retrieval, being exact keyword and boolean matching over typed metadata, answers a query precisely, in that a paper either satisfies the constraints or it does not. Semantic retrieval, being dense vector embeddings that rank documents by meaning, recovers relevance that exact terms miss, but returns what lies nearest in a learned space rather than what provably matches. The two serve different needs. Assembling an exhaustive evidence base, as in a systematic review, demands the precision of the former, while exploring an unfamiliar area rewards the recall of the latter. Deployed literature-search engines force a choice between them.

The aim of this thesis is to develop a scholarly search system that can search based on nested entity queries, and semantic relationships via a three-stage mechanism: 1) Structured Graph Retrieval; 2) Bi-encoder semantic retrieval; 3) Cross-encoder semantic re-ranking. Accordingly, the three research questions that will form an answer to this include:

\begin{enumerate}
    \item To what extent can modern AI construct the artefacts required for deterministic structured graph retrieval of a candidate set of papers: the per-work keyword index extracted offline from the literature, and the intermediate representation formulated from the natural-language query;
    \item How well do current bi-encoders and cross-encoders perform on semantic ad-hoc search tasks on scientific retrieval benchmarks;
    \item How does the three-stage mechanism differ from existing systems, and what explains those differences?
\end{enumerate}

The remainder of this report is organised as follows. Section~\ref{sec:litreview} reviews the literature on scholarly retrieval, positioning the design against dense retrieval, graph-augmented methods, and structured-query translation, and surveying the systems and benchmarks currently available. Section~\ref{sec:research-design} presents the research design: the system architecture, the graph schema and ingestion pipeline, the design of each retrieval stage, and the methodology used to validate them. Section~\ref{sec:data-results} reports the outcome of each component benchmark, which Section~\ref{sec:analysis-discussion} then analyses against the three research questions. Section~VI draws conclusions and identifies future work.

% PLACEHOLDER (conclusion): 'Section~VI' above is left as a literal numeral because
% the Conclusion and Future Work section is not yet in the build. Replace it with
% Section~\ref{sec:conclusion} once that section is written and \input into main.tex. (or \include).
% Assumes an IEEEtran two-column document; contains no preamble.
% Section cross-references use \ref{} to auto-numbered labels. The one exception is
% the Conclusion (Section~VI), left as a literal numeral until that section is written
% and wired into main.tex; see the PLACEHOLDER note at the end of this file.

\section{Introduction}
\label{sec:intro}

A literature search asks two things of a retrieval system: firstly, to find papers that precisely match a set of stated constraints; secondly, to find papers that are about what the user means. These map onto two retrieval paradigms that have proven difficult to reconcile. Structured retrieval, being exact keyword and boolean matching over typed metadata, answers a query precisely, in that a paper either satisfies the constraints or it does not. Semantic retrieval, being dense vector embeddings that rank documents by meaning, recovers relevance that exact terms miss, but returns what lies nearest in a learned space rather than what provably matches. The two serve different needs. Assembling an exhaustive evidence base, as in a systematic review, demands the precision of the former, while exploring an unfamiliar area rewards the recall of the latter. Deployed literature-search engines force a choice between them.

The aim of this thesis is to develop a scholarly search system that can search based on nested entity queries, and semantic relationships via a three-stage mechanism: 1) Structured Graph Retrieval; 2) Bi-encoder semantic retrieval; 3) Cross-encoder semantic re-ranking. Accordingly, the three research questions that will form an answer to this include:

\begin{enumerate}
    \item To what extent can modern AI construct the artefacts required for deterministic structured graph retrieval of a candidate set of papers: the per-work keyword index extracted offline from the literature, and the intermediate representation formulated from the natural-language query;
    \item How well do current bi-encoders and cross-encoders perform on semantic ad-hoc search tasks on scientific retrieval benchmarks;
    \item How does the three-stage mechanism differ from existing systems, and what explains those differences?
\end{enumerate}

The remainder of this report is organised as follows. Section~\ref{sec:litreview} reviews the literature on scholarly retrieval, positioning the design against dense retrieval, graph-augmented methods, and structured-query translation, and surveying the systems and benchmarks currently available. Section~\ref{sec:research-design} presents the research design: the system architecture, the graph schema and ingestion pipeline, the design of each retrieval stage, and the methodology used to validate them. Section~\ref{sec:data-results} reports the outcome of each component benchmark, which Section~\ref{sec:analysis-discussion} then analyses against the three research questions. Section~VI draws conclusions and identifies future work.

% PLACEHOLDER (conclusion): 'Section~VI' above is left as a literal numeral because
% the Conclusion and Future Work section is not yet in the build. Replace it with
% Section~\ref{sec:conclusion} once that section is written and \input into main.tex. (or \include).
% Assumes an IEEEtran two-column document; contains no preamble.
% Section cross-references use \ref{} to auto-numbered labels. The one exception is
% the Conclusion (Section~VI), left as a literal numeral until that section is written
% and wired into main.tex; see the PLACEHOLDER note at the end of this file.

\section{Introduction}
\label{sec:intro}

A literature search asks two things of a retrieval system: firstly, to find papers that precisely match a set of stated constraints; secondly, to find papers that are about what the user means. These map onto two retrieval paradigms that have proven difficult to reconcile. Structured retrieval, being exact keyword and boolean matching over typed metadata, answers a query precisely, in that a paper either satisfies the constraints or it does not. Semantic retrieval, being dense vector embeddings that rank documents by meaning, recovers relevance that exact terms miss, but returns what lies nearest in a learned space rather than what provably matches. The two serve different needs. Assembling an exhaustive evidence base, as in a systematic review, demands the precision of the former, while exploring an unfamiliar area rewards the recall of the latter. Deployed literature-search engines force a choice between them.

The aim of this thesis is to develop a scholarly search system that can search based on nested entity queries, and semantic relationships via a three-stage mechanism: 1) Structured Graph Retrieval; 2) Bi-encoder semantic retrieval; 3) Cross-encoder semantic re-ranking. Accordingly, the three research questions that will form an answer to this include:

\begin{enumerate}
    \item To what extent can modern AI construct the artefacts required for deterministic structured graph retrieval of a candidate set of papers: the per-work keyword index extracted offline from the literature, and the intermediate representation formulated from the natural-language query;
    \item How well do current bi-encoders and cross-encoders perform on semantic ad-hoc search tasks on scientific retrieval benchmarks;
    \item How does the three-stage mechanism differ from existing systems, and what explains those differences?
\end{enumerate}

The remainder of this report is organised as follows. Section~\ref{sec:litreview} reviews the literature on scholarly retrieval, positioning the design against dense retrieval, graph-augmented methods, and structured-query translation, and surveying the systems and benchmarks currently available. Section~\ref{sec:research-design} presents the research design: the system architecture, the graph schema and ingestion pipeline, the design of each retrieval stage, and the methodology used to validate them. Section~\ref{sec:data-results} reports the outcome of each component benchmark, which Section~\ref{sec:analysis-discussion} then analyses against the three research questions. Section~VI draws conclusions and identifies future work.

% PLACEHOLDER (conclusion): 'Section~VI' above is left as a literal numeral because
% the Conclusion and Future Work section is not yet in the build. Replace it with
% Section~\ref{sec:conclusion} once that section is written and \input into main.tex. (or \include).
% Assumes an IEEEtran two-column document; contains no preamble.
% Section cross-references use \ref{} to auto-numbered labels. The one exception is
% the Conclusion (Section~VI), left as a literal numeral until that section is written
% and wired into main.tex; see the PLACEHOLDER note at the end of this file.

\section{Introduction}
\label{sec:intro}

A literature search asks two things of a retrieval system: firstly, to find papers that precisely match a set of stated constraints; secondly, to find papers that are about what the user means. These map onto two retrieval paradigms that have proven difficult to reconcile. Structured retrieval, being exact keyword and boolean matching over typed metadata, answers a query precisely, in that a paper either satisfies the constraints or it does not. Semantic retrieval, being dense vector embeddings that rank documents by meaning, recovers relevance that exact terms miss, but returns what lies nearest in a learned space rather than what provably matches. The two serve different needs. Assembling an exhaustive evidence base, as in a systematic review, demands the precision of the former, while exploring an unfamiliar area rewards the recall of the latter. Deployed literature-search engines force a choice between them.

The aim of this thesis is to develop a scholarly search system that can search based on nested entity queries, and semantic relationships via a three-stage mechanism: 1) Structured Graph Retrieval; 2) Bi-encoder semantic retrieval; 3) Cross-encoder semantic re-ranking. Accordingly, the three research questions that will form an answer to this include:

\begin{enumerate}
    \item To what extent can modern AI construct the artefacts required for deterministic structured graph retrieval of a candidate set of papers: the per-work keyword index extracted offline from the literature, and the intermediate representation formulated from the natural-language query;
    \item How well do current bi-encoders and cross-encoders perform on semantic ad-hoc search tasks on scientific retrieval benchmarks;
    \item How does the three-stage mechanism differ from existing systems, and what explains those differences?
\end{enumerate}

The remainder of this report is organised as follows. Section~\ref{sec:litreview} reviews the literature on scholarly retrieval, positioning the design against dense retrieval, graph-augmented methods, and structured-query translation, and surveying the systems and benchmarks currently available. Section~\ref{sec:research-design} presents the research design: the system architecture, the graph schema and ingestion pipeline, the design of each retrieval stage, and the methodology used to validate them. Section~\ref{sec:data-results} reports the outcome of each component benchmark, which Section~\ref{sec:analysis-discussion} then analyses against the three research questions. Section~VI draws conclusions and identifies future work.

% PLACEHOLDER (conclusion): 'Section~VI' above is left as a literal numeral because
% the Conclusion and Future Work section is not yet in the build. Replace it with
% Section~\ref{sec:conclusion} once that section is written and \input into main.tex.